\section{Consequences for design and analysis}\label{sec:guidance}

\paragraph{When will a readout see the bias?} \Cref{eq:scaling} with $C$ from \cref{tab:arms} answers it. For a leading-edge
or peak time and a lateral spread of 0.5 to 0.8 pitches at full drift, $C$ is 4 to 8, and the compression is 10\,\% at
$N \approx 7$ strips and 1\,\% at $N \approx 25$. A $\mu$TPC with a 5~mm gap and \SI{0.4}{mm} strips has $N = 12\tan\theta$ and
reads 1.3 at \ang{10}, 1.08 at \ang{20} and 1.02 at \ang{40}: the published curve. A pad TPC with a 1~m column and 1~cm pads has
$N = 100\tan\theta$ and a bias below $10^{-3}$ at every angle it measures; the same TPC reading a 3~cm segment near its anode
would see a few per cent. A readout that sums strips to save channels raises $C$ and lowers $N$ at once; in the pitch scan of
\cref{fig:pitch} doubling the pitch raised the end term by 40\,\% at fixed angle.

\paragraph{What the hit times cannot tell you.} The implied drift velocity of the raw ladder is not a drift velocity: on
chamber A it reads 45 and \SI{52}{\micro\metre/ns} for 37 and 36 (\cref{tab:remedies}), and every correction moves it. A
$\mu$TPC that takes $v$ from the ladder will read the right angle at one inclination and the wrong one everywhere else; one
that takes $v$ from the gas tables will read every angle too steep. A calibration against an external angle absorbs the
multiplicative part and leaves the $1/N^2$ part, which no constant can remove. The interior term and a drift-velocity error
have the same signature.

\paragraph{What to record.} The remedies of \cref{sec:remedies} rank the information: a time per strip gives the
compressed ladder; a time and an amplitude allow the edge treatments and the forward model on hit lists (\ang{1.2}); the samples
of the strips around the cluster allow the neighbour subtraction and the full forward fit (\ang{1.0}). Whatever is kept, the
neighbours of the cluster must be kept (the VMM's neighbour-enable logic does this), because the end strips are the
measurement of the spread and the model needs them. A threshold raised to suppress the ghosts removes the information that
locates the column ends.

\paragraph{How to separate the mechanisms.} The angle cannot; the compression is invariant under the exchange of a same-time
footprint for a same-time coupling of equal variance (\cref{sec:mechanisms}). Three measurements can, and each needs only
existing data: the neighbour share against drift depth on near-normal tracks (the constant is the same-time share, the growth is
diffusion, a delay is resistive spreading), the far-neighbour share against the position in the strip (a floor is a fixed
coupling, with its closure against the fitted position), and the channel map (electronics against board). The published
coupling fractions were obtained from none of these; a fraction tuned to restore the angle is a measure of the total same-time
spread, and should be reported as such.

\paragraph{What ``capacitive coupling'' should mean in the next paper.} On chamber A a same-time share of a few per cent per
neighbour is present, its second-neighbour part follows the channel map and is electronic, its first-neighbour part is on the
board and of unidentified origin, and its contribution to the angle bias is one of four of equal standing: strip geometry,
footprint, diffusion and (on the resistive view) RC spreading. A paper that attributes its $\mu$TPC bias to coupling should show
the depth and phase dependences that distinguish it from the other three, or call the whole same-time spread by one name.

\paragraph{What the forward model needs.} A measured impulse response per plane, the three spreads calibrated per detector and
run condition (they trade against each other in the waveform $\chi^2$, so the calibration must be anchored: the depth-resolved
stack fixes the diffusion and the mesh-end share independently of the fit), the drift velocity from the column duration, and
an angle scale from an external reference. With these it is unbiased at the per cent level and limited by diffusion and
ionisation statistics~\cite{paperII}. Without the reference it is still unbiased in shape, which the hits are not.
